problem:  
Let \( X \) be a smooth \(n\)-dimensional Fano manifold with \( [\omega_0] = c_1(X) \), and let \( \omega_t = \omega_0 + i\partial\bar{\partial}\varphi_t \) evolve by the normalized Kähler–Ricci flow  
\[
\frac{\partial \varphi_t}{\partial t} 
= \log\frac{(\omega_0 + i\partial\bar{\partial}\varphi_t)^n}{\omega_0^n}
+ \varphi_t - h_0,
\qquad \int_X e^{-\varphi_t} \omega_0^n = V.
\]
Assume that \( \mathcal{D}(\varphi_t) \), the Ding functional, is bounded below and that the family \(\{\varphi_t\}\) is subconvergent in the \(L^1\)-topology on the finite energy space \(\mathcal{E}^1(X,\omega_0)\) to a weak geodesic ray \(\varphi_\infty(s)\) emanating from \(\omega_0\).  

The Boucksom–Hisamoto–Jonsson correspondence associates to \(\varphi_\infty\) a non‑Archimedean metric \(\phi_\infty^{\mathrm{NA}}\) on \(-K_X\).  

**Question:**  
Is it possible to construct a *formal gradient flow* of the non‑Archimedean Ding functional \(\mathcal{D}^{\mathrm{NA}}\) in the space of non‑Archimedean metrics on \(-K_X\), such that the non‑Archimedean limit \(\phi_\infty^{\mathrm{NA}}\) of the Kähler–Ricci flow arises as a stationary point of this non‑Archimedean flow?  
Equivalently, can one define a natural evolution equation on the space of non‑Archimedean metrics whose stationary points are K‑polystable limits, and whose weak limit coincides with the asymptotic slope direction of the analytic Kähler–Ricci flow?

Why is it a "good" problem:  
This reformulation corrects and sharpens the intuition behind connecting Kähler–Ricci flow dynamics and non‑Archimedean geometry. It removes ill‑defined notions like a full “non‑Archimedean gradient flow continuation” but retains their essential mathematical core: the interplay between gradient flow structures of energy functionals in analytic and non‑Archimedean settings. The problem asks for the existence and characterization of a natural variational evolution—the non‑Archimedean analog of Kähler–Ricci flow—that mirrors the metric–algebraic bridge predicted by stability theory.

If successful, such a theory would thereby formalize the analytic–algebro‑geometric correspondence not only in static (energy-minimizing) terms but also dynamically, offering a non‑Archimedean variational framework mirroring parabolic complex geometric flows. This would open an entirely new line of inquiry joining pluripotential theory, Berkovich analytic geometry, and geometric analysis at a foundational level—a precise yet profoundly forward‑looking research direction.